\section{UXMSS. The Stateful Tree}
Standard XMSS uses a balanced binary tree, which minimizes the maximum path length but imposes a logarithmic cost on every signature. For \texttt{N} signatures, every signature has size proportional to \texttt{log\_2(N)}. SHRINCS utilizes an Unbalanced Merkle Tree (specifically a "Right-Skewed" tree) to optimize for the typical Bitcoin use case: a wallet that performs very few transactions in its lifecycle.

In this topology the \texttt{q}-th signature requires a path of length \texttt{q}. This linear growth is superior to logarithmic growth for small \texttt{q}. For \texttt{q=1}, the path is just 16 bytes. Even for \texttt{q=10}, the path is $10 \times 16 = 160$ bytes, which is extremely small compared to the stateless part.
\subsection{UXMSS Tree Structure}
The UXMSS tree is constructed recursively as a right-skewed binary tree with \texttt{hsf + 1} leaves. To describe the tree structure, we use "level" to denote depth from the root, with level 0 at the root and level \texttt{hsf - 1} at the bottom internal node. Note that this UXMSS convention is the inverse of the convention used in the balanced XMSS trees of the stateless component (Section~7). However, the ADRS height field for internal nodes is set to (\texttt{hsf - level}), so both constructions use increasing height values toward the root.
The UXMSS tree is constructed as follows:
\begin{itemize}
    \item Level 0 is the root public key \texttt{PK.sf}.
    \item Level \texttt{i} (for \texttt{0 <= i <= hsf-2}, internal):
    \begin{itemize}
        \item  Left child: WOTS+C public key for signature \texttt{q=i+1}
        \item  Right child: Root of subtree for signatures \texttt{q>i+1}
    \end{itemize}
    \item Level \texttt{hsf-1} (bottom)
    \begin{itemize}
        \item Left child: WOTS+C public key for signature \texttt{q=hsf}
        \item Right child: WOTS+C public key for signature \texttt{q=hsf+1}
    \end{itemize}
\end{itemize}

This means the tree has \texttt{hsf + 1} leaves (one per valid signature index \texttt{q $\in$ \{1, ..., hsf + 1\}}), and every internal node has a WOTS+C leaf as its left child and a UXMSS subtree (or a second leaf at the bottom) as its right child.

\subsection{UXMSS TreeHash}
% \mknote{Should we have a unified API for these functions with similar logic: xmss treehash(SK.seed, PK.seed, ADRS, target height, start idx). Is level the target hash that is output? If Lvl is 1 what do we output: WOTSpk or rightsubtree.}
The \texttt{uxmss\_treehash} function computes subtree roots within the right-skewed UXMSS tree. Unlike balanced XMSS where treehash recursively computes both subtrees symmetrically, UXMSS treehash follows the right-skewed structure: the left child is always a WOTS+C public key (a leaf), while the right child is either another UXMSS subtree (for internal levels) or a second WOTS+C public key (at the bottom level). The level parameter indicates depth from the root, with level 0 being the root and level \texttt{hsf - 1} being the bottom internal node.

All UXMSS operations use the stateful address types (\texttt{SF\_WOTS\_HASH}, \texttt{SF\_WOTS\_PK}, \texttt{SF\_TREE}) and are fixed at \texttt{layer = 0}, \texttt{tree\_address = 0} (since there is exactly one stateful tree per SHRINCS key pair).

\begin{verbatim}
    Algorithm: uxmss_treehash(SK.seed, PK.seed, ADRS, level)
    Input:
        SK.seed: secret seed
        PK.seed: public seed
        ADRS: tweak (layer = 0, tree_address = 0)
        level: depth from root (0 = root, hsf - 1 = bottom internal node)
    Output:
        node: hash value (n bytes)

    1. left ← wots_PKgen(SK.seed, PK.seed, ADRS, level + 1, SF)
    2. if level == hsf - 1:
            right ← wots_PKgen(SK.seed, PK.seed, ADRS, hsf + 1, SF)
        else:
            right ← uxmss_treehash(SK.seed, PK.seed, ADRS, level + 1)
    3. setTypeAndClear(ADRS, SF_TREE)
    4. setTreeHeight(ADRS, hsf - level)
    5. setTreeIndex(ADRS, 0)
    6. return Tw(PK.seed, ADRS, left || right)
\end{verbatim}
\paragraph{Note:} The keypair index passed to \texttt{wots\_PKgen} corresponds to the signature number \texttt{q}. At level \texttt{i}, the left child is the leaf for \texttt{q = i + 1}. At the bottom level (\texttt{hsf - 1}), the right child is the leaf for \texttt{q = hsf + 1}.

\subsection{UXMSS Root Computation}
The \texttt{uxmss\_root} function computes \texttt{PK.sf}, the root of the stateful UXMSS tree. It initializes the address structure with layer 0 and tree address 0 (since there is only one stateful tree), then invokes \texttt{uxmss\_treehash} starting from level 0. The resulting root, combined with \texttt{PK.sl} from the stateless tree, forms the basis for \texttt{PK.root}.
\begin{verbatim}
    Algorithm: uxmss_root(SK.seed, PK.seed, ADRS)
    Input:
        SK.seed: secret seed
        PK.seed: public seed
        ADRS: tweak
    Output:
        PK.sf: stateful public key

    1. setLayerAddress(ADRS, 0)
    2. setTreeAddress(ADRS, 0)
    3. return uxmss_treehash(SK.seed, PK.seed, ADRS, 0)
\end{verbatim}

\subsection{UXMSS Authentication Path Generation}
For signature number \texttt{1 <= q <= hsf + 1}, the authentication path contains \texttt{min(q, hsf)} sibling nodes needed to reconstruct \texttt{PK.sf} from the WOTS+C public key at position \texttt{q}. The structure of the path differs based on whether \texttt{q <= hsf} or \texttt{q = hsf + 1}. For \texttt{q <= hsf}, the first sibling is either a subtree root (if \texttt{q < hsf}) or the rightmost WOTS+C key (if \texttt{q = hsf}), followed by q-1 WOTS+C public keys ascending toward the root. For \texttt{q = hsf + 1} (the rightmost leaf), all hsf siblings are WOTS+C public keys. This asymmetric structure reflects the right-skewed nature of UXMSS.

\begin{verbatim}
    Algorithm: uxmss_auth_path(SK.seed, PK.seed, ADRS, q)
    Input:
        SK.seed, PK.seed: seeds
        ADRS: address
        q: signature number (1 to hsf + 1)
    Output:
        auth: authentication path (min(q, hsf)*n bytes)

    1. auth ← []
    2. setLayerAddress(ADRS, 0)
    3. setTreeAddress(ADRS, 0)
    4. if q <= hsf:
            if q == hsf:
                auth ← auth || wots_PKgen(SK.seed, PK.seed, ADRS, hsf + 1, SF)
            else:
                auth ← auth || uxmss_treehash(SK.seed, PK.seed, ADRS, q)  
            for i from 1 to q - 1:
                auth ← auth || wots_PKgen(SK.seed, PK.seed, ADRS, q - i, SF)
    5. else:  
            for i from 0 to hsf - 1:
                auth ← auth || wots_PKgen(SK.seed, PK.seed, ADRS, hsf - i, SF)
    6. return auth
\end{verbatim}
\paragraph{Note:} For \texttt{q <= hsf}, the authentication path has exactly \texttt{q} nodes. For \texttt{q = hsf + 1}, it has \texttt{hsf} nodes, since the rightmost leaf shares its parent's other child (a WOTS+C key) rather than a subtree root. Consequently, both \texttt{q = hsf} and \texttt{q = hsf + 1} produce authentication paths of identical byte length (\texttt{hsf * n} bytes). A verifier receiving only the signature cannot determine \texttt{q} from the authentication path length alone when the path has \texttt{hsf} nodes. This ambiguity is resolved during verification by trying both candidate values and accepting whichever reconstructs a valid \texttt{PK.sf} root.

\subsection{UXMSS Public Key Recovery from Signature}
Given a WOTS+C signature and authentication path, this algorithm reconstructs \texttt{PK.sf}. First, it recovers the WOTS+C public key from the signature. Then, starting from this leaf, it iteratively combines the current node with authentication path nodes to compute parent nodes until reaching the root.

The combination order follows from the tree structure:
\begin{itemize}
    \item Case \texttt{q <= hsf}: The recovered key is a left child. The first combination pairs the recovered key (left) with the first auth node (right). All subsequent combinations have the auth node (a WOTS+C key from a lower \texttt{q} value) on the left and the current node on the right.
    \item Case \texttt{q = hsf + 1}: The recovered key is the rightmost leaf (a right child). All auth nodes are WOTS+C keys on the left.
\end{itemize}

\begin{verbatim}
    Algorithm: uxmss_pkFromSig(wots_sig, auth, m, PK.seed, PK.root, ADRS, q)
    Input:
        wots_sig: WOTS+C signature
        auth: authentication path 
        m: signed message
        PK.seed: public seed
        PK.root: root public key
        ADRS: address
        q: signature number
    Output:
        root: computed PK.sf, or "false" if invalid

    1. setLayerAddress(ADRS, 0)
    2. setTreeAddress(ADRS, 0)
    3. pk ← wots_pkFromSig(wots_sig, m, PK.seed, PK.root, ADRS, q, SF, false)
    4. if pk == false:
        return false
    5. node ← pk
    6. setTypeAndClear(ADRS, SF_TREE)
    7. if q <= hsf:
        a. setTreeHeight(ADRS, hsf - (q - 1))
        b. setTreeIndex(ADRS, 0)
        c. node ← Tw(PK.seed, ADRS, node || auth[0 : n])
        for i from 1 to q - 1:
            a. setTreeHeight(ADRS, hsf - (q - 1 - i))
            b. setTreeIndex(ADRS, 0)
            c. node ← Tw(PK.seed, ADRS, auth[i * n : (i + 1) * n] || node)
    8. else:  
        for i from 0 to hsf - 1:
            a. setTreeHeight(ADRS, i + 1)
            b. setTreeIndex(ADRS, 0)
            c. node ← Tw(PK.seed, ADRS, auth[i * n : (i + 1) * n] || node)
    9. return node
\end{verbatim}

\paragraph{Note:} In step 3, \texttt{is\_internal = false} because the stateful UXMSS signs user messages directly (requiring phase 1 + phase 2 hashing). This differs from the stateless XMSS layers which always sign tree roots or FORS public keys (\texttt{is\_internal = true}).

\subsection{UXMSS Signature Generation}
To sign a message using the stateful path, UXMSS signature generation produces a WOTS+C signature followed by the authentication path for position \texttt{q}. The signature number \texttt{q} must be tracked as state and incremented after each use to ensure one-time signature security. As noted in Section~1.2, the state counter should be persisted before generating the signature to prevent key reuse in case of interruption.

The resulting signature size is \texttt{|wots\_sig| + min(q, hsf)*n} bytes, which grows linearly with \texttt{q} up to \texttt{q = hsf}. The last two signatures (\texttt{q = hsf} and \texttt{q = hsf + 1}) have identical size. For the first signature (\texttt{q = 1}), this yields the minimal signature that makes SHRINCS attractive for Bitcoin applications.
\begin{verbatim}
    Algorithm: uxmss_sign(m, SK.seed, SK.prf, PK.seed, PK.root, ADRS, q)
    Input:
        m: message
        SK.seed: secret seed
        SK.prf: PRF key
        PK.seed: public seed
        PK.root: root public key
        ADRS: address
        q: signature number 
    Output:
        sig: UXMSS signature (WOTS+C signature || authentication path)

    1. setLayerAddress(ADRS, 0)
    2. setTreeAddress(ADRS, 0)
    3. wots_sig ← wots_sign(m, SK.seed, SK.prf, PK.seed, PK.root, ADRS, q, SF, false)
    4. auth ← uxmss_auth_path(SK.seed, PK.seed, ADRS, q)
    5. return wots_sig || auth
\end{verbatim}

\subsection{Contribution to Signature Size and Verification Complexity}
Each UXMSS signature consists of a WOTS+C signature plus an authentication path whose length depends on the signature index \texttt{q}:
\begin{center}
    \texttt{total = R\_SIZE + c + l*n + min(q, hsf)*n}
\end{center}

For concrete instances and signatures' number:
\begin{verbatim}
    SHRINCS-B:
        q=1:        32 + 4 + 16*16 + 16 =       308 bytes
        q=10:       32 + 4 + 16*16 + 16*10 =    452 bytes  
        q=158:      32 + 4 + 16*16 + 16*158 =  2820 bytes
        q=159(max): 32 + 4 + 16*16 + 16*158 =  2820 bytes // same as q=158; < 2856 (stateless)
    SHRINCS-L:
        q=1:        32 + 4 + 64*16 + 16 =       1076 bytes
        q=10:       32 + 4 + 64*16 + 16*10 =    1220 bytes  
        q=206:      32 + 4 + 64*16 + 16*206 =   4356 bytes
        q=207(max): 32 + 4 + 64*16 + 16*206 =   4356 bytes // same as q=206; < 4392 (stateless)
\end{verbatim}

Verification cost per UXMSS signature is calculated as WOTS+C verification cost + \texttt{min(q, hsf)} hashes.